\documentclass[a4paper,9.5pt]{article}
\usepackage[margin=1.2cm]{geometry}
\usepackage{fontspec}
\usepackage{polyglossia}
\setmainlanguage{czech}

\usepackage{multicol}
\usepackage{xcolor}
\usepackage{tcolorbox}
\usepackage{enumitem}
\usepackage{fancyhdr}

\setsansfont{DejaVu Sans}
\setmonofont{DejaVu Sans Mono}
\renewcommand{\familydefault}{\sfdefault}

\pagestyle{fancy}
\fancyhf{}
\lhead{\textbf{FRIENDLY FIRE: TAHÁK PRO PŘEDNÁŠEJÍCÍHO} \small(Kryštof Pešek / kof)}
\rhead{\small Zbrojovka, Válečný šrot, Bratrstvo \& Mystery of Memory | 28 slajdů}
\cfoot{\thepage}

\setlength{\columnsep}{0.8cm}
\setlist{nosep,leftmargin=1.1em}

\tcbset{
  colback=black!3,
  colframe=black!70,
  boxrule=0.4pt,
  arc=2pt,
  left=3pt,right=3pt,top=3pt,bottom=3pt,
  before skip=3pt,after skip=4pt
}

\newcommand{\badge}[1]{\textbf{\colorbox{black}{\textcolor{white}{\footnotesize~#1~}}}}

\begin{document}

\begin{multicols}{2}

\begin{tcolorbox}
\badge{00:00 - 02:00} \textbf{Slajd 01: Úvod a rámec}\\
$\bullet$ Ne akademická teorie, ale terénní záchrana válečného stolu.\\
$\bullet$ Motivy: Zbrojovka Brno, Los Alamos, armádní šrot, veteránské bratrstvo a existenciální téma \textit{Mystery of Memory}.
\end{tcolorbox}

\begin{tcolorbox}
\badge{02:00 - 04:00} \textbf{Slajd 02: Krédo mediální archeologie}\\
$\bullet$ \textit{„Dílo neexistuje, nemá-li kontext.“} (23. 7. 2016).\\
$\bullet$ Muzejní hrob (pietní sklo) vs. živý dýchající stroj.\\
$\bullet$ Cílem je obnova původní kinetické a sonické choreografie.
\end{tcolorbox}

\begin{tcolorbox}
\badge{04:00 - 07:00} \textbf{Slajd 03: Zbrojovka Brno \& Otec za války}\\
$\bullet$ Bohuslav narozen 1937 v Brně-Židenicích.\\
$\bullet$ Otec nástrojařem ve Zbrojovce Brno (\textit{Waffenwerke Brünn}) při nucené výrobě pušek Mauser a kulometů pro frontu.\\
$\bullet$ Dětství v náletech, vůně řezných emulzí a špony na montérkách.\\
$\bullet$ Strojní průmyslovka na Sokolské: vyučený strojař. Posedlost vůlemi ozubení a mechanikou předcházela studiu na FAMU.
\end{tcolorbox}

\begin{tcolorbox}
\badge{07:00 - 09:30} \textbf{Slajd 04: The Black Hole v Los Alamos}\\
$\bullet$ Přesun do Nového Mexika (1980): sousedství Los Alamos a Sandia.\\
$\bullet$ Vrakoviště Eda Grothuse: nákup leteckého hliníku z bombardérů B-52, pneumatických ventilů z raket a optických hranolů na váhu.\\
$\bullet$ \textbf{Negativy plošných spojů (PCB photomasks):} Litografické filmy armádních procesorů jako světelné artefakty a archeologie kybernetiky.
\end{tcolorbox}

\begin{tcolorbox}
\badge{09:30 - 12:00} \textbf{Slajd 05: The Brotherhood: 6 stolů \& Epistémický prostor}\\
$\bullet$ Přechod od \textit{single frame} k \textit{epistemic space} a vojenskému haraburdí.\\
$\bullet$ \textbf{6 stolů (1990--1998, ICC Tokio):} T1 Translocations (plotr $17\times 17$, kolejnice), T2 Automata (RPT hlavice), T3 Friendly Fire (dělič, PAD-5), T4 Stealth (wafer stroj), T5 Scribe (rameno, OCR, mizející kód), T6 The Maiden (operační stůl, vějíře, Steina).\\
$\bullet$ \textit{Affectual programming:} přetížení front, ztráta paketů $\rightarrow$ živý organismus.\\
$\bullet$ \textbf{Friendly Fire:} střelba do vlastních řad – zbraň požírající své tvůrce.
\end{tcolorbox}

\begin{tcolorbox}
\badge{12:00 - 14:30} \textbf{Slajd 06: Woody \& The Mystery of Memory}\\
$\bullet$ \textbf{The Art of Memory (1987):} V poušti levitují tělesa s pamětí válek 20. století (španělská válka, Hirošima, atomový hřib).\\
$\bullet$ Paměť jako elektromagnetický prostor a aktivní rekombinace.\\
$\bullet$ Hardware jako vnější protéza mysli pro vytěsněné rány války.
\end{tcolorbox}

\begin{tcolorbox}
\badge{14:30 - 17:00} \textbf{Slajdy 07--08: Anatomie \& 18 let hibernace}\\
$\bullet$ \textbf{Slajd 07:} Hliník z B-52, písty, optika, PCB masky, Linux i386, Roland pad.\\
$\bullet$ \textbf{Slajd 08:} Po ICC Tokio 1998 rozloženo do beden. LaserDisc zničený, Opto 22 spálená, ventily plné prachu. Kurátoři: \textit{„Nemožné zprovoznit.“}
\end{tcolorbox}

\begin{tcolorbox}
\badge{17:00 - 21:00} \textbf{Slajdy 09--12: Expedice \& Tání paměti}\\
$\bullet$ \textbf{Slajd 09:} Telefonát od Jennifer DeFelice přes Miloše Vojtěchovského.\\
$\bullet$ \textbf{Slajd 10 (Tři body):} Rodeo Road (sklad), Ruffina Street (hala), Agua Fria.\\
$\bullet$ \textbf{Slajd 11 (Zlom): Woody v roce 2016 ztrácí paměť (demence).} Tvůrce strojové paměti sleduje rozpad vlastních vzpomínek. Pamatuje si strojařinu, ale konfigurace a cesty se rozpadají v mlze.\\
$\bullet$ \textbf{Slajd 12 (Steina):} Lidský index skladů. Bez ní by se nenašla jediná čočka.
\end{tcolorbox}

\columnbreak

\begin{tcolorbox}
\badge{21:00 - 26:00} \textbf{Slajdy 13--17: Forenzní archeologie \& Disk}\\
$\bullet$ \textbf{Slajd 13:} Nález plotnového disku i386 v bedně se zdroji. Obraz přes \texttt{ddrescue}.\\
$\bullet$ \textbf{Slajd 14 (Rekonstrukce):} Když mozek zapomíná, mluví magnetická stopa. Skládání z Woodyho záblesků, Steininy sluchové paměti a bitů.\\
$\bullet$ \textbf{Slajd 15:} Emulace v QEMU přímo v hale na laptopu.\\
$\bullet$ \textbf{Slajd 16 (UNICOM v C):} \texttt{logicdd} (:9850 blackboard), \texttt{mididd} (:9800 Roland pad), \texttt{serialdd} (:9900 LaserDisc), \texttt{optodd} (\texttt{/dev/lp2}, bit 6 napájení, \textbackslash232 video, \textbackslash254 diaky, \textbackslash247 karusel), \texttt{picdd}, dynamičtí \texttt{actors} s \texttt{PID kill}.\\
$\bullet$ \textbf{Slajd 17:} Rozeta \texttt{PITT.TXT} (8. 3. 1997). Časový kód 922 sekund (15m 22s).
\end{tcolorbox}

\begin{tcolorbox}
\badge{26:00 - 31:00} \textbf{Slajdy 18--21: Etika \& Náhrady}\\
$\bullet$ \textbf{Slajd 18 (Doktrína):} \textit{Zachovat tělo; nahradit zaniklé orgány otevřeným systémem.}\\
$\bullet$ \textbf{Slajd 19 (Opto 22 $\rightarrow$ Arduino + SSR):} Charles Tuttle navrhl miniaturní desku s polovodičovými relé. Konec proudových rázů, přechod z LPT na USB serial.\\
$\bullet$ \textbf{Slajd 20 (LaserDisc $\rightarrow$ RPi H.264):} Nález NTSC DV kazety; deinterlace \texttt{yadif}; \texttt{omxplayer} s D-Bus seekem (<80 ms).\\
$\bullet$ \textbf{Slajd 21:} Kódový diff. Zastaralé proprietární C vs. čistý POSIX shell a Arduino.
\end{tcolorbox}

\begin{tcolorbox}
\badge{31:00 - 38:00} \textbf{Slajdy 22--24: Zásah autorů, Video \& Audio}\\
$\bullet$ \textbf{Slajd 22 [VIDEO MONTÁŽ]:} Spustit 65s video sestřih montáže v dílně.\\
$\bullet$ \textbf{Slajd 23 (Steinino ucho):} Příkaz ztlumit drone a vytáhnout dialogy vojáků.\\
$\bullet$ \textbf{Slajd 24 [AUDIO UKÁZKA]:} Pustit 60s ukázku z rozhovoru na Zoom H2n (\texttt{vasulka\_interview\_excerpt.opus}). Woodyho inženýrský humor a vzpomínky na Brno.
\end{tcolorbox}

\begin{tcolorbox}
\badge{38:00 - 44:00} \textbf{Slajdy 25--27: Špilberk, Light Revisited \& Divák}\\
$\bullet$ \textbf{Slajd 25:} Transport 67 kg beden přes oceán (\texttt{crates.txt}).\\
$\bullet$ \textbf{Slajd 26:} Špilberk 2016 (\textit{Mystery of Memory}). Autonomie 8 hod/den (\texttt{run.sh}). Instalace \textbf{Light Revisited – Noisefields}: prostorový rozklad šumu \textbf{dichroickými zrcadly (1/3 spektra RGB)}.\\
$\bullet$ \textbf{Slajd 27:} Roland MIDI pad. Divák udeří $\rightarrow$ písty syčí, clony padají, video skáče. Divák se stává součástí zbraňového aparátu.
\end{tcolorbox}

\begin{tcolorbox}
\badge{44:00 - 50:00} \textbf{Slajd 28: Epilog \& Islandský odkaz}\\
$\bullet$ Woody Vašulka zemřel v prosinci 2019 v Santa Fe.\\
$\bullet$ \textbf{Island:} \textit{Vasulka Chamber} v Reykjavíku; výstava a sympozium \textbf{BERG Contemporary (2024--2026)} s Tate a LACMA.\\
$\bullet$ Paměť tvůrce odešla, ale dílo žije dál díky otevřenému kódu (\texttt{github.com/k0f/friendlyfire}) a technické genealogii (Joey Heinen). Otevření diskuse.
\end{tcolorbox}

\begin{tcolorbox}[colback=black!8]
\textbf{Pohotovostní historky a záludné dotazy:}\\
1. \textit{„Proč negativy plošných spojů?“} $\rightarrow$ Posvátné svitky – zhmotnění logických hradel dřív, než jsou vyleptána do křemíku.\\
2. \textit{„Jak se lišil od militaristů?“} $\rightarrow$ Jako syn dělníka ze Zbrojovky znal dřinu v kovu. \textit{The Brotherhood} je chladnou pitvou válečné dehumanizace.\\
3. \textit{„Proč ne LaserDisc z eBay?“} $\rightarrow$ Laserová dioda po 20 letech degraduje. RPi s H.264 nabízí 50 let bezúdržbového chodu.\\
4. \textit{„Zrcadla v Light Revisited?“} $\rightarrow$ Dichroická zrcadla dělí spektrum na třetiny (R, G, B) $\rightarrow$ 3 prostorově oddělené kruhy šumu.\\
5. \textit{„Role Joeyho Heinena a Islandu?“} $\rightarrow$ Heinen (NYU 2014) sepsal genealogii UNICOMu; BERG Contemporary v Reykjavíku 2024--2026 uspořádala velkou záchrannou výstavu s Tate/LACMA.
\end{tcolorbox}

\end{multicols}

\vspace{0.1cm}
\centering\footnotesize\textcolor{gray}{--τέχνη \& kof | Vygenerováno pro Kryštofa Peška}

\end{document}
